%\subsubsection{Verifica dei documenti}
%Ogni documento prodotto dal team passa per una o più fasi di verifica prima di essere pubblicato ufficialmente nel repository \textit{GitHub}.\\
%Tale fasi vengono realizzate solamente in seguito alla consegna del documento da parte dei responsabili di stesura, i quali ritengono che in quel momento il prodotto da loro realizzato sia completato.\\
%La verifica dei documenti prevede che venga assicurato:
%\begin{itemize}
%    \item Il rispetto delle norme tipografiche stabilite nelle \textit{Norme di progetto} del team;
%    \item Il rispetto della struttura definita nelle \textit{Norme di progetto} del team;
%    \item Il rispetto del contenuto, il quale deve essere chiaro, coerente, preciso ed esaustivo;
%    \item Il rispetto della sintassi italiana e della correttezza ortografica dei termini utilizzati.
%\end{itemize}

\documentclass[10pt, a4paper]{article}

% Parametri che modificano il file main.tex
% Le uniche parti da cambiare su main.tex sono:
% - tabella versioni
% - testo file

\def\titolo{Analisi dei \\ \ requisiti} % \\ per andare a capo, \ per spazio
\def\titoloHeader{Analisi dei requisiti}
\def\data{v2.0.0(0)}

% \def\listaComponenti{
% Bresolin G.,
% Campese M.,
% Ciriolo I.,
% Dugo A.,
% Feltrin E.,
% Michelon R.,
% Orlandi G.
% }

\setcounter{secnumdepth}{5}
\setcounter{tocdepth}{5}
\makeatletter
\newcommand\subsubsubsection{\@startsection{paragraph}{4}{\z@}{-2.5ex\@plus -1ex \@minus -.25ex}{1.25ex \@plus .25ex}{\normalfont\normalsize\bfseries}}
\newcommand\subsubsubsubsection{\@startsection{subparagraph}{5}{\z@}{-2.5ex\@plus -1ex \@minus -.25ex}{1.25ex \@plus .25ex}{\normalfont\normalsize\bfseries}}
\makeatother

\usepackage{style}
\usepackage{headerfooter}
\usepackage{comment}

\title{\titolo}
\author{SWEetCode}

\begin{document}

% PRIMA PAGINA
\begin{titlepage}
    \thispagestyle{empty}
    \begin{tikzpicture}[remember picture, overlay]
        % TRIANGOLI
        \draw[fill=secondarycolor, secondarycolor] (current page.north west) -- (current page.south west) -- (8.8, -28);
        \draw[fill=primarycolor, primarycolor] (-3, 5) -- (4, -13.6) -- (11, 5);

        % LOGO
        \node [xshift=-5cm, yshift=25cm] (logo) at (current page.south east) {\includegraphics[width=6.5cm]{img/logo.png}};

        % SWEETCODE - DATE
        \node [anchor=north east, align=right, xshift=-1.2cm, yshift=20.5cm, text=black] (sweetcode) at (current page.south east) {\fontsize{32pt}{36pt}\selectfont SWEetCode};
        \draw[line width=4pt, lightcol] ([xshift=-3cm, yshift=-0.37cm]sweetcode.south west) -- ([yshift=-0.37cm]sweetcode.south east);
        \node [anchor=north east, align=right, xshift=-1.2cm, yshift=18.7cm, text=black] (date) at (current page.south east){\fontsize{24pt}{24pt} \selectfont Verbale \tipoVerb};

        % NOME FILE
        \node [anchor=north east, text width=15cm, align=right, xshift=-1.2cm, yshift=17cm, text=black] (titolo) at (current page.south east){\fontsize{48pt}{48pt}\textbf{\data}};

        % BOX DATI PARTECIPANTI
        \ifthenelse{\equal{\tipoVerb}{Esterno}}{
            \node[anchor=north east, xshift=-1.2cm, yshift=14.5cm, minimum width=8cm] (box) at (current page.south east){};
        }{
            \node[anchor=north east, xshift=-1.2cm, yshift=12.5cm, minimum width=8cm] (box) at (current page.south east){};
        }

        % RESPONSABILE
        \node[anchor=north west, align=left] (dati2) at (box.north west) {\fontsize{15pt}{15pt}\selectfont \textbf{Responsabile}};
        \draw[line width=4pt, lightcol] (dati2.south west) -- ([xshift=8cm]dati2.south west);
        \node[anchor=north west, align=left] (dati21) at (dati2.south west){\fontsize{13pt}{13pt}\selectfont \nomeResp};

        % VERIFICATORE
        \node[anchor=north west, yshift=-1cm, align=left] (dati3) at (dati21.north west) {\fontsize{15pt}{15pt}\selectfont \textbf{Verificatore}};
        \draw[line width=4pt, lightcol] (dati3.south west) -- ([xshift=8cm]dati3.south west);
        \node[anchor=north west, align=left] (dati31) at (dati3.south west){\fontsize{13pt}{13pt}\selectfont \nomeVer};

        % SEGRETARIO DI RIUNIONE
        \node[anchor=north west, yshift=-1cm, align=left] (dati4) at (dati31.north west) {\fontsize{15pt}{15pt}\selectfont \textbf{Segretario di Riunione}};
        \draw[line width=4pt, lightcol] (dati4.south west) -- ([xshift=8cm]dati4.south west);
        \node[anchor=north west, align=left] (dati41) at (dati4.south west){\fontsize{13pt}{13pt}\selectfont \nomeSegr};
        
        % UNIPD - SWE
        \node [xshift=4.4cm, yshift=2.3cm, draw, secondarycolor, text=white] (uni) at (current page.south west) {\fontsize{20pt}{20pt} \selectfont Università di Padova};
        \node [xshift=0.65cm, yshift=0.7cm, draw, secondarycolor, text=white, below=of uni] (corso) {\fontsize{20pt}{20pt}\selectfont Ingegneria del Software};

        % FIRMA AZIENDA
        \ifthenelse{\equal{\tipoVerb}{Esterno}}{
            \draw[line width=4pt, lightcol] ([xshift=-1.2cm, yshift=7cm]current page.south east) -- ([xshift=-8cm, yshift=7cm]current page.south east);
            \node [xshift=-4.8cm, yshift=8cm] (logo) at (current page.south east) {\includegraphics[width=6cm]{img/firme/\firmaAzienda}};
            \node[anchor=north west, xshift=12.9cm, yshift=6.65cm, align=left] at (current page.south west)
            {\fontsize{13pt}{13pt}\selectfont L'azienda: \nomeAzienda};
        }

        % FIRMA CODING COWBOYS
        \draw[line width=4pt, lightcol] ([xshift=-1.2cm, yshift=4.35cm]current page.south east) -- ([xshift=-8cm, yshift=4.35cm]current page.south east);
        \node [xshift=-4.8cm, yshift=5.5cm] (logo) at (current page.south east) {\includegraphics[width=6cm]{img/firme/menon.png}};
        \node[anchor=north west, xshift=12.9cm, yshift=4cm, align=left] at (current page.south west)
        {\fontsize{13pt}{13pt}\selectfont Coding Cowboys: Menon G.};

        % FIRMA
        \draw[line width=4pt, lightcol] ([xshift=-1.2cm, yshift=1.8cm]current page.south east) -- ([xshift=-8cm, yshift=1.8cm]current page.south east);
        \node [xshift=-4.8cm, yshift=2.45cm] (logo) at (current page.south east) {\includegraphics[width=6cm]{img/firme/\firmaResp}};
        \node[anchor=north west, xshift=12.9cm, yshift=1.45cm, align=left] at (current page.south west)
        {\fontsize{13pt}{13pt}\selectfont Il Responsabile: \nomeResp};
        
    \end{tikzpicture}
\end{titlepage}

% REGISTRO DELLE VERSIONI
%Registro in ordine dalla più recente alla meno recente!
{\renewcommand{\arraystretch}{1.5}
\section*{Registro delle versioni}

\begin{xltabular}{\textwidth}{c|c|c|c|X}
\label{tab:long}

\textbf{Versione} & \textbf{Data} & \quantities{\textbf{Responsabile di}\\\textbf{stesura}}& \textbf{Revisore} & \quantities{\textbf{Dettaglio e}\\\textbf{motivazioni}} \\
\endfirsthead

\textbf{Versione} & \textbf{Data} & \quantities{\textbf{Responsabile di}\\\textbf{stesura}}& \textbf{Revisore} & \quantities{\textbf{Dettaglio e}\\\textbf{motivazioni}} \\
\endhead

\multicolumn{5}{r}{{Continua nella pagina successiva}} \\
\endfoot

\endlastfoot

\hline
v1.10.0(7) & $2023-12-29$ & \quantities{Michelon R.} & Dugo A. & Aggiunti diagrammi use case e di attività.\\
\hline
v1.10.0(6) & $2023-12-29$ & \quantities{Michelon R.} & Dugo A. & Aggiornamento use case e requisiti.\\
\hline
v1.0.0(14) & $2023-12-06$ & \quantities{Bresolin G.} & Orlandi G. & Aggiornamento UC e requisiti presenti. Introduzione diagrammi dei casi d'uso. Aggiunta di nuovi UC e nuovi requisiti.\\
\hline
v1.0.0(5) & $2023-11-23$ & \quantities{Ciriolo I.} & Campese M. & Introduzione, Descrizione, Studio dei primi Casi d'uso, Studio dei primi requisiti.\\

    
\end{xltabular}}
\newpage

% INDICE
\tableofcontents
\newpage
\listoffigures
\newpage
\listoftables
\newpage

% INTRODUZIONE
\section{Introduzione}
\subsection{Scopo del documento}
L'obiettivo che ci si pone nella realizzazione di questo documento è la definizione delle metriche di valutazione e validazione del progetto, e la specifica degli obiettivi di qualità del prodotto finale. Tali parametri vengono stabiliti in accordo ai requisiti e 
alle aspettative del proponente e talvolta a discrezione del team sulla base delle valutazioni fatte nel corso di studi.

\subsection{Glossario}
Al fine di evitare incomprensioni relative alla terminologia usata all’interno del documento viene fornito un Glossario nel file apposito. La prima occorrenza di gni termine presente in tale documento presenta una scrittura in corsivo ed un pedice |g|.

\subsection{Riferimenti}
   \subsubsection{Riferimenti normativi}
   \begin{itemize}
    \item \textit{Norme di progetto};
    \item \textit{Regolamento del progetto didattico}: \\
    \href{https://www.math.unipd.it/~tullio/IS-1/2023/Dispense/PD2.pdf}{https://www.math.unipd.it/~tullio/IS-1/2023/Dispense/PD2.pdf};
    \item \textit{Guidelines for
    the application of ISO 9001:2000 to
    computer software}: \\
    \href{https://cdn.standards.iteh.ai/samples/35867/36860aa4caba4c84b26051db576456d3/ISO-IEC-90003-2004.pdf}{https://cdn.standards.iteh.ai/samples/35867/36860aa4caba4c84b26051db576456d3/ISO-IEC-90003-2004.pdf}.
    \end{itemize}
    
    \subsubsection{Riferimenti informativi}
    \begin{itemize}
    \item \textit{Presentazione capitolato C1}: \\
    \href{https://www.math.unipd.it/~tullio/IS-1/2023/Progetto/C1p.pdf}{https://www.math.unipd.it/~tullio/IS-1/2023/Progetto/C1p.pdf};
    \item \textit{Analisi dei requisiti};
    \item \textit{Piano di progetto};
    \item \textit{Verbali esterni ed interni};
    \item \textit{Qualità di processo (argomento T8)}: \\
    \href{https://www.math.unipd.it/~tullio/IS-1/2023/Dispense/T8.pdf}{https://www.math.unipd.it/~tullio/IS-1/2023/Dispense/T8.pdf};
    \item \textit{Verifica e Validazione: Introduzione (argomento T9)}: \\
    \href{https://www.math.unipd.it/~tullio/IS-1/2023/Dispense/T9.pdf}{https://www.math.unipd.it/~tullio/IS-1/2023/Dispense/T9.pdf}.
    
   
    \end{itemize}

    
% Obiettivi di qualità
\newpage
\section{Obiettivi metrici di qualità}
\label{ObiettiviQualità}
Per garantire la qualità di processo si è deciso di aderire allo standard ISO/IEC 90003:2004. In questa sezione si presentano i valori accettabili e ideali delle metriche che si rifanno a questo standard; le definizioni di tali metriche sono riportate nella sezione apposita (\S Lista delle metriche) del documento Norme di progetto(\color{red}Aggiungere versione\color{black}).
\subsection{Qualità di processo}
%\paragraph{}breve descrizione qualità di processo da inserire.\\
 {\renewcommand{\arraystretch}{1.5}
 \begin{table}[H]
    \begin{tabularx}{\textwidth}{p{0.1\textwidth}|p{0.5\textwidth}|X|X}
    \textbf{ID} & \textbf{Nome metrica} & \textbf{Valore tollerabile} & \textbf{Valore ambito}  \\
    \hline
     M1PMS & Percentuale di metriche soddisfatte & $ \ge75\% $& $\ge95\% $\\
    \hline
     M2RNP & Rischi non previsti & $ \le2 $& $\le1 $\\
    \hline
     M3VP & Variazione di piano & $ \ge-8 $ & $ \ge0 $ \\
    \hline
     M4VC & Variazione di costo & $0$ & $\le0 $ \\
    \hline
     M5ISR &  Indice di stabilità dei requisiti & $ \ge70\% $ & $ 100\% $  \\
    \hline
     M6CCM & Complessità ciclomatica media & $\le6 $ & $\le3 $ \\
    \hline
    M7SC & Statement coverage & $ \ge75\% $ & $ 100\% $ \\
    
    \hline
     M8BC & Branch Coverage & $ \ge75\% $ & $ 100\% $ \\
    \end{tabularx}
    \caption{Metriche per la qualità di processo}
    \end{table}
    
\subsection{Qualità di prodotto}
Abbiamo analizzato quali caratteristiche fossero necessarie per la realizzazione di un prodotto di 
qualità.\\
 {\renewcommand{\arraystretch}{1.5}
 \begin{table}[H]
  \begin{tabularx}{\textwidth}{p{0.1\textwidth}|p{0.5\textwidth}|X|X}
    \textbf{ID} & \textbf{Nome metrica} & \textbf{Valore tollerabile} & \textbf{Valore ambito}  \\
    %\endfirsthead
    \hline
     M9IG & Indice di Gulpease & $ \ge40\% $& $\ge80\% $\\
    \hline
     M10PRI & Percentuale di requisiti implementati & $ \ge85\% $ & $ 100\% $ \\
    \hline
     M11LMC & Linee medie di codice per metodo & $\le30$ & $\le20$ \\
    \hline
    M12DC & Densità dei commenti & $\ge25\%$ & $\ge30\%$ \\
    \hline
M13TUS & Percentuale di test superati & $ \ge85\% $& $100\%$\\

    \end{tabularx}
    \caption{Metriche per la qualità di prodotto}
    \end{table}

    
%\subsection{Software}





\subsection{Qualità per obiettivo}
\subsubsection{Processi primari}

\subsubsubsection{Acquisizione}
 
\subsubsubsection{Fornitura}


{\renewcommand{\arraystretch}{1.5}
\begin{table}[H]
\begin{xltabular}{\textwidth}{p{0.1\textwidth}|p{0.5\textwidth}|X|X}


\textbf{ID} & \textbf{Nome metrica} & \textbf{Valore tollerabile} & \textbf{Valore ambito}   \\
\endfirsthead

\textbf{ID} & \textbf{Nome metrica} & \textbf{Valore tollerabile} & \textbf{Valore ambito}   \\
\endhead

\multicolumn{4}{r}{{Continua nella pagina successiva}} \\
\endfoot

\endlastfoot

\hline
 M2RNP & Rischi non previsti & $ \le2 $& $\le1 $\\
\hline
 M3VP & Variazione di piano & $ \ge-8 $ & $ \ge0 $ \\
\hline
 M4VC & Variazione di costo & $0$ & $\le0 $ \\

\end{xltabular}
\caption{Metriche per la fornitura}
\end{table}


 
\subsubsubsection{Sviluppo}

\subsubsubsubsection{Analisi dei requisiti}
{\renewcommand{\arraystretch}{1.5}
\begin{table}[H]
\begin{xltabular}{\textwidth}{p{0.1\textwidth}|p{0.5\textwidth}|X|X}

\textbf{ID} & \textbf{Nome metrica} & \textbf{Valore tollerabile} & \textbf{Valore ambito}   \\
\endfirsthead

\textbf{ID} & \textbf{Nome metrica} & \textbf{Valore tollerabile} & \textbf{Valore ambito}   \\
\endhead

\multicolumn{4}{r}{{Continua nella pagina successiva}} \\
\endfoot

\endlastfoot
  \hline
 M5ISR &  Indice di stabilità dei requisiti & $ \ge70\% $ & $ 100\% $  \\

\end{xltabular}
\caption{Metriche per l'analisi dei requisiti}
\end{table}


\subsubsubsubsection{Progettazione}

{\renewcommand{\arraystretch}{1.5}
\begin{table}[H]
\begin{xltabular}{\textwidth}{p{0.1\textwidth}|p{0.5\textwidth}|X|X}

\textbf{ID} & \textbf{Nome metrica} & \textbf{Valore tollerabile} & \textbf{Valore ambito}   \\
\endfirsthead

\textbf{ID} & \textbf{Nome metrica} & \textbf{Valore tollerabile} & \textbf{Valore ambito}   \\
\endhead

\multicolumn{4}{r}{{Continua nella pagina successiva}} \\
\endfoot

\endlastfoot
\hline
M10PRI & Percentuale di requisiti implementati & $ \ge85\% $ & $ 100\% $ \\ 
\end{xltabular}
\caption{Metriche per la progettazione}
\end{table}

\subsubsubsubsection{Codifica}
{\renewcommand{\arraystretch}{1.5}
\begin{table}[H]
\begin{xltabular}{\textwidth}{p{0.1\textwidth}|p{0.5\textwidth}|X|X}

\textbf{ID} & \textbf{Nome metrica} & \textbf{Valore tollerabile} & \textbf{Valore ambito}   \\
\endfirsthead

\textbf{ID} & \textbf{Nome metrica} & \textbf{Valore tollerabile} & \textbf{Valore ambito}   \\
\endhead

\multicolumn{4}{r}{{Continua nella pagina successiva}} \\
\endfoot

\endlastfoot
\hline
 M6CCM & Complessità ciclomatica media & $\le6 $ & $\le3 $ \\
\hline
M7SC & Statement coverage & $ \ge75\% $ & $ 100\% $ \\  
\hline
M8BC & Branch Coverage & $ \ge75\% $ & $ 100\% $ \\
\hline
M10PRI & Percentuale di requisiti implementati & $ \ge85\% $ & $ 100\% $ \\ \hline
M11LMC & Linee medie di codice per metodo & $\le30$ & $\le20$ \\
\hline
M12DC & Densità dei commenti & $\ge25\%$ & $\ge30\%$ \\ 
\end{xltabular}
\caption{Metriche per la codifica}
\end{table}

\subsubsubsubsection{Testing}
{\renewcommand{\arraystretch}{1.5}
\begin{table}[H]
\begin{xltabular}{\textwidth}{p{0.1\textwidth}|p{0.5\textwidth}|X|X}

\textbf{ID} & \textbf{Nome metrica} & \textbf{Valore tollerabile} & \textbf{Valore ambito}   \\
\endfirsthead

\textbf{ID} & \textbf{Nome metrica} & \textbf{Valore tollerabile} & \textbf{Valore ambito}   \\
\endhead

\multicolumn{4}{r}{{Continua nella pagina successiva}} \\
\endfoot

\endlastfoot

\hline
M13TUS & Percentuale di test & $ \ge85\% $& $100\%$\\

    
\end{xltabular}
\caption{Metriche per il testing}
\end{table}
 
    {\renewcommand{\arraystretch}{1.5}
    \begin{tabularx}{\textwidth}{p{0.05\textwidth}|p{0.35\textwidth}|X|X}
    \textbf{ID} & \textbf{Nome metrica} & \textbf{Valore accettabile} & \textbf{Valore ottimale}  \\
    \hline
    \multicolumn{4}{l}{\cellcolor{primarycolor}\textbf{\textit{Fornitura}}} \\
    \hline
     &  &  &  \\
    \hline
     &  &  &  \\
    \hline
     &  &  &  \\
    \hline
     &  &  &  \\
    \hline
     &  &  &  \\
    \hline
     &  &  &  \\
    \hline
     &  &  &  \\
    \hline
    \multicolumn{4}{l}{\cellcolor{primarycolor}\textbf{\textit{Sviluppo}}} \\
    \hline
     &  &  &  \\
    \hline
     &  &  &  \\
    \end{tabularx}}

% PROCESSI DI SUPPORTO
\subsubsection{Processi di supporto}
\subsubsubsection{Documentazione}
{\renewcommand{\arraystretch}{1.5}
\begin{table}[H]
\begin{xltabular}{\textwidth}{p{0.1\textwidth}|p{0.5\textwidth}|X|X}

\textbf{ID} & \textbf{Nome metrica} & \textbf{Valore tollerabile} & \textbf{Valore ambito}   \\
\endfirsthead

\textbf{ID} & \textbf{Nome metrica} & \textbf{Valore tollerabile} & \textbf{Valore ambito}   \\
\endhead

\multicolumn{4}{r}{{Continua nella pagina successiva}} \\
\endfoot

\endlastfoot
\hline
M9IG & Indice di Gulpease & $ \ge40\% $& $\ge80\% $\\
\end{xltabular}
\caption{Metriche per la documentazione}
\end{table}

\begin{comment}
%\subsubsection{Documentazione}
    {\renewcommand{\arraystretch}{1.5}
    \begin{tabularx}{\textwidth}{p{0.18\textwidth}|p{0.6\textwidth}|X}
    \textbf{Obiettivo} & \textbf{Descrizione} & \textbf{Metriche}  \\
    \hline
    Correttezza Risposte &   & \\
    \hline
    Velocità di esecuzione &  &  \\
    \end{tabularx}}
    {\renewcommand{\arraystretch}{1.5}
    \begin{tabularx}{\textwidth}{p{0.18\textwidth}|p{0.6\textwidth}|X}
    \textbf{Obiettivo} & \textbf{Descrizione} & \textbf{Metriche}  \\
    \hline
    Verifica &  &  \\
    \hline
    Gestione della qualità &  &  \\
    \end{tabularx}}

\end{comment}    
 % PROCESSI ORGANIZZATIVI   
\subsubsection{Processi organizzativi}
{\renewcommand{\arraystretch}{1.5}
\begin{table}[H]
\begin{xltabular}{\textwidth}{p{0.1\textwidth}|p{0.5\textwidth}|X|X}

\textbf{ID} & \textbf{Nome metrica} & \textbf{Valore tollerabile} & \textbf{Valore ambito}   \\
\endfirsthead

\textbf{ID} & \textbf{Nome metrica} & \textbf{Valore tollerabile} & \textbf{Valore ambito}   \\
\endhead

\multicolumn{4}{r}{{Continua nella pagina successiva}} \\
\endfoot

\endlastfoot
\hline
 M2RNP & Rischi non previsti & $ \le2 $& $\le1 $\\
\end{xltabular}
\caption{Metriche per i processi organizzativi }
\end{table}
\begin{comment}
    {\renewcommand{\arraystretch}{1.5}
    \begin{tabularx}{\textwidth}{p{0.18\textwidth}|p{0.6\textwidth}|X}
    \textbf{Obiettivo} & \textbf{Descrizione} & \textbf{Metriche}  \\
    \hline
    Gestione organizzativa &  &  \\
    \end{tabularx}}
    \end{comment}
    


    

\begin{comment}
    {\renewcommand{\arraystretch}{1.5}
    \begin{tabularx}{\textwidth}{p{0.05\textwidth}|p{0.35\textwidth}|X|X}
    \textbf{ID} & \textbf{Nome metrica} & \textbf{Valore accettabile} & \textbf{Valore ottimale}  \\
    \hline
    \multicolumn{4}{l}{\cellcolor{primarycolor}\textbf{\textit{Gestione organizzativa}}} \\
    \hline
     &  &  &  \\
    \end{tabularx}}  
\end{comment}




%Test e specifiche
\newpage
\section{Verifica}


\subsection{Test di unità}
Questa sezione vuota verrà ampliata in futuro con la progettazione di dettaglio del prodotto.
\begin{comment}
    

    {\renewcommand{\arraystretch}{1.5}
    \begin{tabularx}{\textwidth}{p{0.12\textwidth}|p{0.7\textwidth}|X}
    \textbf{ID} & \textbf{Descrizione} & \textbf{Stato}  \\
    \hline
     &  & \\
    \hline
     &  &  \\
    \hline
     &  & \\
    \hline
     &  &  \\
    \end{tabularx}}
    \subsubsection{Tracciamento}
\end{comment}    
\subsection{Test di integrazione}
Questa sezione vuota verrà ampliata in futuro con la progettazione di dettaglio del prodotto.
\begin{comment}
    {\renewcommand{\arraystretch}{1.5}
    \begin{tabularx}{\textwidth}{p{0.12\textwidth}|p{0.7\textwidth}|X}
    \textbf{ID} & \textbf{Descrizione} & \textbf{Stato}  \\
    \hline
     &  & \\
    \hline
     &  &  \\
    \hline
     &  & \\
    \end{tabularx}}
    \subsubsection{Tracciamento}
\end{comment}








\subsection{Test di sistema}
Il documento Analisi dei Requisiti identifica i requisiti che devono essere completamente coperti dai test di sistema. Di seguito è presente un elenco dei test di sistema:
    \begin{xltabular}{\textwidth}{|c|X|c|c|}
    %\begin{xltabular}{\textwidth}{p{0.12\textwidth}|p{0.55\textwidth}|p{0.12\textwidth}|X}
    \caption{Tabella dei test di sistema}
    \label{tab:test_sistema}\\
    \hline
    \textbf{ID} & \textbf{Descrizione} & \textbf{Requisito} & \textbf{Stato}  \\
    \hline
    \endfirsthead
    \caption[]{Tabella dei test di sistema (cont)}\\
    \hline
    \textbf{ID} & \textbf{Descrizione} & \textbf{Requisito} & \textbf{Stato}  \\
    \hline
    \endhead
    \multicolumn{4}{r}{\footnotesize\textit{Continua nella pagina successiva}}
    \endfoot
    \hline
    \endlastfoot
    TS.1 & L’utente deve poter configurare i componenti del sistema (modello LLM, vector database, modello di embeddings, sistema di archiviazione) al primo avvio  dell’applicazione. & RF.O.1 & I \\
    \hline
    TS.2 & L’utente deve poter configurare il modello LLM. & RF.O.2 & I \\
    \hline
    TS.3 & L’utente deve poter visualizzare la lista dei modelli LLM disponibili. & RF.O.2.1 & I \\
    \hline
    TS.4 & L’utente deve poter visualizzare i singoli modelli LLM presenti in lista. & RF.O.2.1.1 & I \\
    \hline
    TS.5 & L’utente deve poter scegliere un modello LLM tra quelli disponibili. & RF.O.2.2 & I \\
    \hline
    TS.6 & L’utente deve poter inserire le credenziali e i parametri richiesti dal modello LLM. & RF.O.2.3 & I \\
    \hline
    TS.7 & L’utente deve poter inserire le credenziali e i parametri richiesti dal modello LLM di OpenAI. & RF.O.2.3.1 & I \\
    \hline
    TS.8 & L’utente deve poter inserire l’API key di un account OpenAI. & RF.O.2.3.1.1 & I \\
    \hline
    TS.9 & L’utente deve poter inserire il nome del modello di LLM di OpenAI. & RF.O.2.3.1.2 & I \\
    \hline
    TS.10 & L’utente deve poter inserire le credenziali e i parametri richiesti dal modello LLM di HuggingFace. & RF.O.2.3.2 & I \\
    \hline
    TS.11 & L’utente deve poter inserire l’API key di un account HuggingFace. & RF.O.2.3.2.1 & I \\
    \hline
    TS.12 & L’utente deve poter inserire il nome del modello di LLM di HuggingFace. & RF.O.2.3.2.2 & I \\
    \hline
    TS.13 & L’utente deve poter inserire le credenziali e i parametri richiesti dal modello LLM locale. & RF.O.2.3.3 & I \\
    \hline
    TS.14 & L’utente deve poter confermare la scelta del modello LLM. & RF.O.2.4 & I \\
    \hline
    TS.15 & Il tentativo di configurazione del modello LLM deve fallire se l’utente ha inserito credenziali o parametri non validi. & RF.O.3 & I \\
    \hline
    TS.16 & L’utente deve visualizzare un errore nel caso in cui abbia inserito parametri o credenziali non validi del modello LLM. & RF.O.3.1 & I \\
    \hline
    TS.17 & L’utente deve poter configurare il vector database. & RF.O.4 & I \\
    \hline
    TS.18 & L’utente deve poter visualizzare la lista dei vector database disponibili. & RF.O.4.1 & I \\
    \hline
    TS.19 & L’utente deve poter visualizzare i singoli vector database presenti in lista. & RF.O.4.1.1 & I \\
    \hline
    TS.20 & L’utente deve poter scegliere un vector database tra quelli disponibili. & RF.O.4.2 & I \\
    \hline
    TS.21 & L’utente deve poter inserire le credenziali e i parametri richiesti dal vector database. & RF.O.4.3 & I \\
    \hline
    TS.22 & L’utente deve poter inserire le credenziali e i parametri richiesti dal vector database Pinecone. & RF.O.4.3.1 & I \\
    \hline
    TS.23 & L’utente deve poter inserire l’API key di un account Pinecone. & RF.O.4.3.1.1 & I \\
    \hline
    TS.24 & L’utente deve poter inserire il nome dell’index del database di Pinecone. & RF.O.4.3.1.2 & I \\
    \hline
    TS.25 & L’utente deve poter inserire il nome dell’environment di Pinecone. & RF.O.4.3.1.3 & I \\
    \hline
    TS.26 & L’utente deve poter inserire database ChromaDB. & RF.O.4.3.2 & I \\
    \hline
    TS.27 & L’utente deve poter inserire il percorso della cartella di persistenza per il database di ChromaDB. & RF.O.4.3.2.1 & I \\
    \hline
    TS.28 & L’utente deve poter inserire il nome della collection di ChromaDB. & RF.O.4.3.2.2 & I \\
    \hline
    TS.29 & L’utente deve poter confermare la scelta del vector database. & RF.O.4.4 & I \\
    \hline
    TS.30 & Il tentativo di configurazione del vector database deve fallire se l’utente ha inserito credenziali o parametri non validi. & RF.O.5 & I \\
    \hline
    TS.31 & L’utente deve visualizzare un errore nel caso in cui abbia inserito parametri o credenziali non validi del vector database. & RF.O.5.1 & I \\
    \hline
    TS.32 & L’utente deve poter configurare il modello di embeddings. & RF.O.6 & I \\
    \hline
    TS.33 & L’utente deve poter visualizzare la lista dei modelli di embeddings disponibili. & RF.O.6.1 & I \\
    \hline
    TS.34 & L’utente deve poter visualizzare i singoli modelli di embeddings presenti in lista. & RF.O.6.1.1 & I \\
    \hline
    TS.35 & L’utente deve poter scegliere un modello di embeddings tra quelli disponibili. & RF.O.6.2 & I \\
    \hline
    TS.36 & L’utente deve poter inserire le credenziali e i parametri richiesti dal modello di embeddings. & RF.O.6.3 & I \\
    \hline
    TS.37 & L’utente deve poter inserire le credenziali e i parametri richiesti dal modello di embeddings di OpenAI. & RF.O.6.3.1 & I \\
    \hline
    TS.38 & L’utente deve poter inserire l’API key di un account OpenAI. & RF.O.6.3.1.1 & I \\
    \hline
    TS.39 & L’utente deve poter inserire il nome del modello di embeddings di OpenAI. & RF.O.6.3.1.2 & I \\
    \hline
    TS.40 & L’utente deve poter inserire le credenziali e i parametri richiesti dal modello di embeddings di HuggingFace. & RF.O.6.3.2 & I \\
    \hline
    TS.41 & L’utente deve poter inserire l’API key di un account HuggingFace. & RF.O.6.3.2.1 & I \\
    \hline
    TS.42 & L’utente deve poter inserire il nome del modello di embeddings di HuggingFace. & RF.O.6.3.2.2 & I \\
    \hline
    TS.43 & L’utente deve poter confermare la scelta del modello di embeddings. & RF.O.6.4 & I \\
    \hline
    TS.44 & Il tentativo di configurazione del modello di embeddings deve fallire se l’utente ha inserito credenziali o parametri non validi. & RF.O.7 & I \\
    \hline
    TS.45 & L’utente deve visualizzare un errore nel caso in cui abbia inserito parametri o credenziali non validi del modello di embeddings. & RF.O.7.1 & I \\
    \hline
    TS.46 & L’utente deve poter configurare la scelta del sistema di archiviazione documenti. & RF.O.8 & I \\
    \hline
    TS.47 & L’utente deve poter visualizzare la lista dei sistemi di archiviazione disponibili. & RF.O.8.1 & I \\
    \hline
    TS.48 & L’utente deve poter visualizzare i singoli sistemi di archiviazione presenti in lista. & RF.O.8.1.1 & I \\
    \hline
    TS.49 & L’utente deve poter scegliere un sistema di archiviazione tra quelli disponibili. & RF.O.8.2 & I \\
    \hline
    TS.50 & L’utente deve poter inserire le credenziali e i parametri richiesti dal sistema di archiviazione. & RF.O.8.3 & I \\
    \hline
    TS.51 & L’utente deve poter inserire le credenziali e i parametri richiesti dal sistema di archiviazione AWS S3. & RF.O.8.3.1 & I \\
    \hline
    TS.52 & L’utente deve poter inserire l’access key id di AWS S3. & RF.O.8.3.1.1 & I \\
    \hline
    TS.53 & L’utente deve poter inserire la secret access key di AWS S3. & RF.O.8.3.1.2 & I \\
    \hline
    TS.54 & L’utente deve poter inserire il nome del bucket di AWS S3. & RF.O.8.3.1.3 & I \\
    \hline
    TS.55 & L’utente deve poter inserire il nome della regione di AWS S3. & RF.O.8.3.1.4 & I \\
    \hline
    TS.56 & L’utente deve poter inserire le credenziali e i parametri richiesti dal sistema di archiviazione locale. & RF.O.8.3.2 & I \\
    \hline
    TS.57 & L’utente deve poter confermare la scelta del sistema di archiviazione. & RF.O.8.4 & I \\
    \hline
    TS.58 & Il tentativo di configurazione del sistema di archiviazione deve fallire se l’utente ha inserito credenziali o parametri non validi. & RF.O.9 & I \\
    \hline
    TS.59 & L’utente deve visualizzare un errore nel caso in cui abbia inserito parametri o credenziali non validi del sistema di archiviazione. & RF.O.9.1 & I \\
    \hline
    TS.60 & L’utente deve poter cambiare la scelta del modello LLM. & RF.O.10 & I \\
    \hline
    TS.61 & L’utente deve poter cambiare la scelta del vector database. & RF.O.11 & I \\
    \hline
    TS.62 & L’utente deve poter cambiare la scelta del modello di embeddings. & RF.O.12 & I \\
    \hline
    TS.63 & L’utente deve poter cambiare la scelta del sistema di archiviazione documenti. & RF.O.13 & I \\
    \hline
    TS.64 & L’utente deve poter visualizzare le caratteristiche di un modello LLM. & RF.O.14 & I \\
    \hline
    TS.65 & L’utente deve poter visualizzare il nome di un modello LLM. & RF.O.14.1 & I \\
    \hline
    TS.66 & L’utente deve poter visualizzare l’organizzazione proprietaria di un modello LLM. & RF.O.14.2 & I \\
    \hline
    TS.67 & L’utente deve poter visualizzare la descrizione di un modello LLM. & RF.O.14.3 & I \\
    \hline
    TS.68 & L’utente deve poter visualizzare la tipologia di un modello LLM. & RF.O.14.4 & I \\
    \hline
    TS.69 & L’utente deve poter visualizzare l’indicatore di costi di un modello LLM. & RF.O.14.5 & I \\
    \hline
    TS.70 & L’utente deve poter visualizzare le caratteristiche di un vector database. & RF.O.15 & I \\
    \hline
    TS.71 & L’utente deve poter visualizzare il nome di un vector database. & RF.O.15.1 & I \\
    \hline
    TS.72 & L’utente deve poter visualizzare l’organizzazione proprietaria di un vector database. & RF.O.15.2 & I \\
    \hline
    TS.73 & L’utente deve poter visualizzare la descrizione di un vector database. & RF.O.15.3 & I \\
    \hline
    TS.74 & L’utente deve poter visualizzare la tipologia di un vector database. & RF.O.15.4 & I \\
    \hline
    TS.75 & L’utente deve poter visualizzare l’indicatore di costi di un vector database. & RF.O.15.5 & I \\
    \hline
    TS.76 & L’utente deve poter visualizzare le caratteristiche di un modello di embeddings. & RF.O.16 & I \\
    \hline
    TS.77 & L’utente deve poter visualizzare il nome di un modello di embeddings. & RF.O.16.1 & I \\
    \hline
    TS.78 & L’utente deve poter visualizzare l’organizzazione proprietaria di un modello di embeddings. & RF.O.16.2 & I \\
    \hline
    TS.79 & L’utente deve poter visualizzare la descrizione di un modello di embeddings. & RF.O.16.3 & I \\
    \hline
    TS.80 & L’utente deve poter visualizzare la tipologia di un modello di embeddings. & RF.O.16.4 & I \\
    \hline
    TS.81 & L’utente deve poter visualizzare l’indicatore di costi di un modello di embeddings. & RF.O.16.5 & I \\
    \hline
    TS.82 & L’utente deve poter visualizzare le caratteristiche di un sistema di archiviazione. & RF.O.17 & I \\
    \hline
    TS.83 & L’utente deve poter visualizzare il nome di un sistema di archiviazione. & RF.O.17.1 & I \\
    \hline
    TS.84 & L’utente deve poter visualizzare l’organizzazione proprietaria di un sistema di archiviazione. & RF.O.17.2 & I \\
    \hline
    TS.85 & L’utente deve poter visualizzare la descrizione di un sistema di archiviazione. & RF.O.17.3 & I \\
    \hline
    TS.86 & L’utente deve poter visualizzare la tipologia di un sistema di archiviazione. & RF.O.17.4 & I \\
    \hline
    TS.87 & L’utente deve poter visualizzare l’indicatore di costi di un sistema di archiviazione. & RF.O.17.5 & I \\
    \hline
    TS.88 & L’utente deve poter visualizzare la configurazione attuale del sistema. & RF.O.18 & I \\
    \hline
    TS.89 & L’utente deve poter visualizzare la configurazione attuale del modello LLM. & RF.O.18.1 & I \\
    \hline
    TS.90 & L’utente deve poter visualizzare la configurazione attuale del vector database. & RF.O.18.2 & I \\
    \hline
    TS.91 & L’utente deve poter visualizzare la configurazione attuale del modello di embeddings. & RF.O.18.3 & I \\
    \hline
    TS.92 & L’utente deve poter visualizzare la configurazione attuale del sistema di archiviazione documenti. & RF.O.18.4 & I \\
    \hline
    TS.93 & L’utente deve poter visualizzare la lista dei documenti presenti nel sistema di archiviazione. & RF.O.19 & I \\
    \hline
    TS.94 & L’utente deve poter visualizzare i singoli documenti presenti nel sistema di archiviazione. & RF.O.19.1 & I \\
    \hline
    TS.95 & L’utente deve poter visualizzare il titolo di un documento presente nel sistema di archiviazione. & RF.O.19.1.1 & I \\
    \hline
    TS.96 & L’utente deve poter visualizzare le dimensioni di un documento presente nel sistema di archiviazione. & RF.O.19.1.2 & I \\
    \hline
    TS.97 & L’utente deve poter visualizzare la data di caricamento di un documento presente nel sistema di archiviazione. & RF.O.19.1.3 & I \\
    \hline
    TS.98 & L’utente deve poter visualizzare se un documento è stato occultato. & RF.O.19.1.4 & I \\
    \hline
    TS.99 & L’utente deve poter caricare un documento nell’area di staging. & RF.O.20 & I \\
    \hline
    TS.100 & L’utente deve poter caricare nell’area di staging un documento di tipo .PDF. & RF.O.20.1 & I \\
    \hline
    TS.101 & L’utente deve poter caricare nell’area di staging un documento di tipo .docx. & RF.O.20.2 & I \\
    \hline
    TS.102 & Il tentativo di caricamento in area di staging di un documento deve fallire se il formato di questo non è supportato. & RF.O.21 & I \\
    \hline
    TS.103 & L’utente deve visualizzare un errore nel caso in cui abbia provato a caricare nell’area di staging un documento con formato non supportato. & RF.O.21.1 & I \\
    \hline
    TS.104 & L’utente deve poter rimuovere un documento caricato nell’area di staging. & RF.O.22 & I \\
    \hline
    TS.105 & L’utente deve poter fare l’upload nel sistema di archiviazione e nel vector database dei documenti caricati nell’area di staging. & RF.O.23 & I \\
    \hline
    TS.106 & L’utente deve poter scegliere se sostituire con il documento caricato un documento con lo stesso nome già presente nel sistema. & RF.O.23.1 & I \\
    \hline
    TS.107 & L’utente deve poter visualizzare il contenuto di ciascun documento presente nel sistema. & RF.O.24 & I \\
    \hline
    TS.108 & L’utente deve poter eliminare un documento presente nel sistema. & RF.O.25 & I \\
    \hline
    TS.109 & L’utente deve poter occultare un documento presente nel sistema. & RF.O.26 & I \\
    \hline
    TS.110 & L’utente deve poter riabilitare un documento presente nel sistema. & RF.O.27 & I \\
    \hline
    TS.111 & L’utente deve poter ricercare un documento presente nel sistema. & RF.O.28 & I \\
    \hline
    TS.112 & L’utente deve poter visualizzare la lista delle chat presenti nell’applicazione. & RF.O.29 & I \\
    \hline
    TS.113 & L’utente deve poter visualizzare le singole chat presenti in lista. & RF.O.29.1 & I \\
    \hline
    TS.114 & L’utente deve poter visualizzare il titolo di una chat. & RF.O.29.1.1 & I \\
    \hline
    TS.115 & L’utente deve poter visualizzare il timestamp dell’ultimo messaggio di una chat. & RF.O.29.1.2 & I \\
    \hline
    TS.116 & L’utente deve poter creare una nuova chat. & RF.O.30 & I \\
    \hline
    TS.117 & Il tentativo di creazione di una nuova chat deve fallire se esiste già una chat vuota. & RF.O.31 & I \\
    \hline
    TS.118 & L’utente deve visualizzare un errore nel caso in cui abbia provato a creare una nuova chat in presenza di una chat vuota. & RF.O.31.1 & I \\
    \hline
    TS.119 & L’utente deve poter eliminare una chat. & RF.O.32 & I \\
    \hline
    TS.120 & L’utente deve poter rinominare una chat. & RF.O.33 & I \\
    \hline
    TS.121 & Il tentativo di rinominazione di una chat deve fallire se il titolo inserito appartiene già ad un’altra chat. & RF.O.34 & I \\
    \hline
    TS.122 & L’utente deve visualizzare un errore nel caso in cui abbia provato a rinominare una chat con un titolo che appartiene già ad un’altra chat. & RF.O.34.1 & I \\
    \hline
    TS.123 & Il tentativo di rinominazione di una chat deve fallire se il titolo inserito è vuoto. & RF.O.35 & I \\
    \hline
    TS.124 & L’utente deve visualizzare un errore nel caso in cui abbia provato a rinominare una chat con un titolo vuoto. & RF.O.35.1 & I \\
    \hline
    TS.125 & L’utente deve poter ricercare una chat e il suo contenuto. & RF.O.36 & I \\
    \hline
    TS.126 & L’utente deve poter visualizzare il contenuto di una chat. & RF.O.37 & I \\
    \hline
    TS.127 & L’utente deve poter visualizzare i singoli messaggi di una chat. & RF.O.37.1 & I \\
    \hline
    TS.128 & L’utente deve poter visualizzare il contenuto testuale di un messaggio presente in una chat. & RF.O.37.1.1 & I \\
    \hline
    TS.129 & L’utente deve poter visualizzare l’orario di un messaggio presente in una chat. & RF.O.37.1.2 & I \\
    \hline
    TS.130 & La risposta fornita dal chatbot deve presentare un riferimento ai documenti da cui ha estrapolato le informazioni esposte. & RF.D.38 & I \\
    \hline
    TS.131 & L’utente deve poter visualizzare le domande suggerite in una chat vuota. & RF.D.39 & I \\
    \hline
    TS.132 & L’utente deve poter inviare una domanda suggerita in una chat vuota. & RF.D.40 & I \\
    \hline
    TS.133 & L’utente deve poter scegliere una delle domande suggerite in una chat vuota. & RF.D.40.1 & I \\
    \hline
    TS.134 & L’utente deve poter confermare l’invio di una domanda suggerita in una chat vuota. & RF.D.40.2 & I \\
    \hline
    TS.135 & L’utente deve poter ascoltare tramite audio un messaggio. & RF.D.41 & I \\
    \hline
    TS.136 & L’utente deve poter copiare un messaggio. & RF.D.42 & I \\
    \hline
    TS.137 & L’utente deve poter attivare la modalità vocale. & RF.OP.43 & I \\
    \hline
    TS.138 & L’utente deve poter disattivare la modalità vocale. & RF.OP.44 & I \\
    \hline
    TS.139 & L’utente deve poter inviare un messaggio in una chat. & RF.O.45 & I \\
    \hline
    TS.140 & L’utente deve poter inserire un messaggio in una chat. & RF.O.45.1 & I \\
    \hline
    TS.141 & L’utente deve poter inserire un messaggio testuale in una chat. & RF.O.45.1.1 & I \\
    \hline
    TS.142 & L’utente deve poter inserire un messaggio in una chat tramite comando vocale. & RF.O.45.1.2 & I \\
    \hline
    TS.143 & L’utente deve poter confermare l’invio del messaggio inserito. & RF.O.45.2 & I \\
    \end{xltabular}
    }

    
    \subsubsection{Tracciamento}
    {\renewcommand{\arraystretch}{1.5}
    \begin{tabularx}{\textwidth}{p{0.12\textwidth}|X}
    \textbf{ID Test} & \textbf{Requisito}  \\
    \hline
    TS.1 & RF.O.1 \\
    \hline
    TS.2 & RF.O.2 \\
    \hline
    TS.3 & RF.O.2.1 \\
    \hline
    TS.4 & RF.O.2.1.1 \\
    \hline
    TS.5 & RF.O.2.2 \\
    \hline
    TS.6 & RF.O.2.3 \\
    \hline
    TS.7 & RF.O.2.3.1 \\
    \hline
    TS.8 & RF.O.2.3.1.1 \\
    \hline
    TS.9 & RF.O.2.3.1.2 \\
    \hline
    TS.10 & RF.O.2.3.2 \\
    \hline
    TS.11 & RF.O.2.3.2.1 \\
    \hline
    TS.12 & RF.O.2.3.2.2 \\
    \hline
    TS.13 & RF.O.2.3.3 \\
    \hline
    TS.14 & RF.O.2.4 \\
    \hline
    TS.15 & RF.O.3 \\
    \hline
    TS.16 & RF.O.3.1 \\
    \hline
    TS.17 & RF.O.4 \\
    \hline
    TS.18 & RF.O.4.1 \\
    \hline
    TS.19 & RF.O.4.1.1 \\
    \hline
    TS.20 & RF.O.4.2 \\
    \hline
    TS.21 & RF.O.4.3 \\
    \hline
    TS.22 & RF.O.4.3.1 \\
    \hline
    TS.23 & RF.O.4.3.1.1 \\
    \hline
    TS.24 & RF.O.4.3.1.2 \\
    \hline
    TS.25 & RF.O.4.3.1.3 \\
    \hline
    TS.26 & RF.O.4.3.2 \\
    \hline
    TS.27 & RF.O.4.3.2.1 \\
    \hline
    TS.28 & RF.O.4.3.2.2 \\
    \hline
    TS.29 & RF.O.4.4 \\
    \hline
    TS.30 & RF.O.5 \\
    \hline
    TS.31 & RF.O.5.1 \\
    \hline
    TS.32 & RF.O.6 \\
    \hline
    TS.33 & RF.O.6.1 \\
    \hline
    TS.34 & RF.O.6.1.1 \\
    \hline
    TS.35 & RF.O.6.2 \\
    \hline
    TS.36 & RF.O.6.3 \\
    \hline
    TS.37 & RF.O.6.3.1 \\
    \hline
    TS.38 & RF.O.6.3.1.1 \\
    \hline
    TS.39 & RF.O.6.3.1.2 \\
    \hline
    TS.40 & RF.O.6.3.2 \\
    \hline
    TS.41 & RF.O.6.3.2.1 \\
    \hline
    TS.42 & RF.O.6.3.2.2 \\
    \hline
    TS.43 & RF.O.6.4 \\
    \hline
    TS.44 & RF.O.7 \\
    \hline
    TS.45 & RF.O.7.1 \\
    \hline
    TS.46 & RF.O.8 \\
    \hline
    TS.47 & RF.O.8.1 \\
    \hline
    TS.48 & RF.O.8.1.1 \\
    \hline
    TS.49 & RF.O.8.2 \\
    \hline
    TS.50 & RF.O.8.3 \\
    \hline
    TS.51 & RF.O.8.3.1 \\
    \hline
    TS.52 & RF.O.8.3.1.1 \\
    \hline
    TS.53 & RF.O.8.3.1.2 \\
    \hline
    TS.54 & RF.O.8.3.1.3 \\
    \hline
    TS.55 & RF.O.8.3.2 \\
    \hline
    TS.56 & RF.O.8.3.2.1 \\
    \hline
    TS.57 & RF.O.8.4 \\
    \hline
    TS.58 & RF.O.9 \\
    \hline
    TS.59 & RF.O.9.1 \\
    \hline
    TS.60 & RF.O.10 \\
    \hline
    TS.61 & RF.O.11 \\
    \hline
    TS.62 & RF.O.12 \\
    \hline
    TS.63 & RF.O.13 \\
    \hline
    TS.64 & RF.O.14 \\
    \hline
    TS.65 & RF.O.14.1 \\
    \hline
    TS.66 & RF.O.14.2 \\
    \hline
    TS.67 & RF.O.14.3 \\
    \hline
    TS.68 & RF.O.14.4 \\
    \hline
    TS.69 & RF.O.14.5 \\
    \hline
    TS.70 & RF.O.15 \\
    \hline
    TS.71 & RF.O.15.1 \\
    \hline
    TS.72 & RF.O.15.2 \\
    \hline
    TS.73 & RF.O.15.3 \\
    \hline
    TS.74 & RF.O.15.4 \\
    \hline
    TS.75 & RF.O.15.5 \\
    \hline
    TS.76 & RF.O.16 \\
    \hline
    TS.77 & RF.O.16.1 \\
    \hline
    TS.78 & RF.O.16.2 \\
    \hline
    TS.79 & RF.O.16.3 \\
    \hline
    TS.80 & RF.O.16.4 \\
    \hline
    TS.81 & RF.O.16.5 \\
    \hline
    TS.82 & RF.O.17 \\
    \hline
    TS.83 & RF.O.17.1 \\
    \hline
    TS.84 & RF.O.17.2 \\
    \hline
    TS.85 & RF.O.17.3 \\
    \hline
    TS.86 & RF.O.17.4 \\
    \hline
    TS.87 & RF.O.17.5 \\
    \hline
    TS.88 & RF.O.18 \\
    \hline
    TS.89 & RF.O.18.1 \\
    \hline
    TS.90 & RF.O.18.2 \\
    \hline
    TS.91 & RF.O.18.3 \\
    \hline
    TS.92 & RF.O.18.4 \\
    \hline
    TS.93 & RF.O.19 \\
    \hline
    TS.94 & RF.O.19.1 \\
    \hline
    TS.95 & RF.O.19.1.1 \\
    \hline
    TS.96 & RF.O.19.1.2 \\
    \hline
    TS.97 & RF.O.19.1.3 \\
    \hline
    TS.98 & RF.O.19.1.4 \\
    \hline
    TS.99 & RF.O.20 \\
    \hline
    TS.100 & RF.O.20.1 \\
    \hline
    TS.101 & RF.O.20.2 \\
    \hline
    TS.102 & RF.O.21 \\
    \hline
    TS.103 & RF.O.21.1 \\
    \hline
    TS.104 & RF.O.22 \\
    \hline
    TS.105 & RF.O.23 \\
    \hline
    TS.106 & RF.O.23.1 \\
    \hline
    TS.107 & RF.O.24 \\
    \hline
    TS.108 & RF.O.25 \\
    \hline
    TS.109 & RF.O.26 \\
    \hline
    TS.110 & RF.O.27 \\
    \hline
    TS.111 & RF.O.28 \\
    \hline
    TS.112 & RF.O.29 \\
    \hline
        TS.113 & RF.O.29 \\
    \hline
    TS.114 & RF.O.29.1 \\
    \hline
    TS.115 & RF.O.29.1.1 \\
    \hline
    TS.116 & RF.O.29.1.2 \\
    \hline
    TS.117 & RF.O.30 \\
    \hline
    TS.118 & RF.O.31 \\
    \hline
    TS.119 & RF.O.31.1 \\
    \hline
    TS.120 & RF.O.32 \\
    \hline
    TS.121 & RF.O.33 \\
    \hline
    TS.122 & RF.O.34 \\
    \hline
    TS.123 & RF.O.34.1 \\
    \hline
    TS.124 & RF.O.35 \\
    \hline
    TS.125 & RF.O.35.1 \\
    \hline
    TS.126 & RF.O.36 \\
    \hline
    TS.127 & RF.O.37 \\
    \hline
    TS.128 & RF.O.37.1 \\
    \hline
    TS.129 & RF.O.37.1.1 \\
    \hline
    TS.130 & RF.O.37.1.2 \\
    \hline
    TS.131 & RF.D.38 \\
    \hline
    TS.132 & RF.D.39 \\
    \hline
    TS.133 & RF.D.40 \\
    \hline
    TS.134 & RF.D.40.1 \\
    \hline
    TS.135 & RF.D.40.2 \\
    \hline
    TS.136 & RF.D.41 \\
    \hline
    TS.137 & RF.D.42 \\
    \hline
    TS.138 & RF.OP.43 \\
    \hline
    TS.139 & RF.OP.44 \\
    \hline
    TS.140 & RF.O.45 \\
    \hline
    TS.141 & RF.O.45.1 \\
    \hline
    TS.142 & RF.O.45.1.1 \\
    \hline
    TS.143 & RF.O.45.1.2 \\
    \end{tabularx}
    }


\subsection{Test di accettazione}
    {\renewcommand{\arraystretch}{1.5}
    \begin{tabularx}{\textwidth}{p{0.12\textwidth}|p{0.7\textwidth}|X}
    \textbf{ID} & \textbf{Descrizione} & \textbf{Stato}  \\
    \hline
    TA.1 & Il sistema deve offrire opzioni che garantiscono la privacy delle informazioni sensibili condivise dall’utente per l’addestramento del modello. & I \\
    \hline
    TA.2 & Il sistema deve garantire il servizio per almeno un utente. & I \\
    \hline
    TA.3 & In caso di caduta dalle connessione, l’applicazione va in fallback su un modello AI locale. & I \\
    \hline
    TA.4 & Nel caso in cui l’utente formuli una richiesta non pertinente al dataset di addestramento, il chatbot deve fornire una risposta cortese in cui espone la sua incapacità nel fornire una risposta alla domanda. & I \\
    \hline
    TA.5 & Il chatbot deve essere in grado di fare caching: se l’utente formula più volte la stessa domanda, il chatbot fornisce la risposta elaborata alla prima occorrenza della domanda. & I \\
    \hline
    TA.6 & L’eliminazione di un documento dallo spazio di storage deve portare all’eliminazione dei metadati del relativo documento, ma non all’eliminazione degli embeddings generati dallo stesso. & I \\
    \hline
    TA.7 & Il prodotto software deve rispettare tutto ciò che è descritto da (Norme di progetto). & I \\
    \hline
    TA.8 & Il prodotto software deve rispettare tutto ciò che è descritto da (Piano di qualifica). & I \\
    \hline
    TA.9 & L’applicazione web deve garantire la presenza di almeno un modello AI locale. & I \\
    \hline
    TA.10 & Le chat devono essere salvate in automatico in un database. & I \\
    \hline
    TA.11 & I documenti caricati devono essere memorizzati in un sistema di archiviazione. & I \\
    \hline
    TA.12 & Il sistema deve garantire il servizio per almeno un utente. & I \\
    \end{tabularx}}

    \subsubsection{Tracciamento}
    {\renewcommand{\arraystretch}{1.5}
    \begin{tabularx}{\textwidth}{p{0.12\textwidth}|X}
    \textbf{ID Test} & \textbf{Requisito}  \\
    \hline
    TA.1 & RQ.O.1 \\
    \hline
    TA.2 & RQ.O.2 \\
    \hline
    TA.3 & RQ.OP.3 \\
    \hline
    TA.4 & RQ.D.4 \\
    \hline
    TA.5 & RQ.OP.5 \\
    \hline
    TA.6 & RQ.D.6 \\
    \hline
    TA.7 & RV.O.1 \\
    \hline
    TA.8 & RV.O.2 \\
    \hline
    TA.9 & RV.O.3 \\
    \hline
    TA.10 & RV.O.4 \\
    \hline
    TA.11 & RV.O.5 \\
    \hline
    TA.12 & RP.O.1 \\
    \end{tabularx}}


\subsection{Checklist }
Questa sezione contiene le liste di controllo usate per l'ispezione dei documenti e del codice; questa è una tecnica di analisi statica che permette di esporre in maniera sistematica gli errori più comuni eventualmente presenti nel prodotto in esame.
\subsubsection{Struttura documentazione}

{\renewcommand{\arraystretch}{1.5}
\begin{table}[H]
\begin{xltabular}{\textwidth}{c|X}


\textbf{Titolo} & \textbf{Descrizione}   \\
\endfirsthead

\textbf{Titolo} & \textbf{Descrizione}   \\
\endhead

\multicolumn{2}{r}{{Continua nella pagina successiva}} \\
\endfoot

\endlastfoot

\hline
\textbf{Riferimenti a documenti} &  Quando viene menzionato un documento versionabile del progetto, al suo titolo va affiancato il numero di versione a cui si riferisce.\\
\hline
\textbf{Riferimenti a risorse web} &  Quando viene menzionato una risorsa web, presumibilmente soggetta a cambiamenti frequenti, a questa va affiancata la data di ultimo accesso alla risorsa stessa.\\
\hline
\textbf{Metriche non ammissibili} & Questa circostanza si può verificare quando si ha l’uscita dal range di accettazione degli indici di misurazione descritti alla sezione \S \hyperref[ObiettiviQualità]{Obiettivi di qualità}; In particolare riferendosi all'indice di Gulpease.\\
\hline
\textbf{Didascalie mancanti} & Tutte le immagini e tabelle devono avere la propria didascalia.\\
\hline
\textbf{Vuoti documentativi} & Non devono essere presenti documenti che presentano sezioni vuote.\\

\hline
\textbf{Numerazione pagine/slide} & Ogni documento o presentazione deve presentare a piè di pagina una numerazione scritta come segue \textit{"\# pagina corrente / \# pagine totali"}.\\
\hline
\textbf{Occorrenze multiple con pedice} & Quando un termine appare più volte in un documento e ha un pedice, il pedice deve essere presente solamente nella prima occorrenza.\\
\hline
\textbf{Punteggiatura elenchi} & Gli elenchi devono terminare con un punto se completano la frase introduttiva, altrimenti con una virgola o punto e virgola.\\
\hline
\textbf{Formato delle date errato} & Tutte le date presenti nella documentazione prevedono il seguente formato: \textit{"AAAA-MM-GG"}.\\
\hline
\textbf{Menzioni a documenti errate} & Tutte le menzioni a documenti devono avere la lettera iniziale del documento in maiuscolo e le restanti lettere in minuscolo.\\
\hline
\textbf{Citazione scorretta di persona }& Ogni menzione di una persona deve essere formulata come segue: \textit{"Cognome N."} \\

\hline
\textbf{Ordine alfabetico} & (\textbf{da controllare}) Gli elenchi di nomi e le tabelle devono essere ordinati in ordine alfabetico a meno che non sia specificato diversamente o le circostanze richiedono altre tipologie di ordinamento.\\
\hline
\textbf{Ordine task per scadenza} & I task devono essere ordinati in base alla loro scadenza, dal più vicino al più lontano.\\

\end{xltabular}
\caption{Punti di controllo per l'ispezione della struttura dei documenti}
\end{table}


\subsubsection{Errori di forma e ortografici}

{\renewcommand{\arraystretch}{1.5}
\begin{table}[H]
\begin{xltabular}{\textwidth}{c|X}


\textbf{Titolo} & \textbf{Descrizione}   \\
\endfirsthead

\textbf{Titolo} & \textbf{Descrizione}   \\
\endhead

\multicolumn{2}{r}{{Continua nella pagina successiva}} \\
\endfoot

\endlastfoot


\hline
\textbf{Repo sos. maschile} & Il termine "repository" deve essere usato nel genere maschile.\\
\hline
\textbf{Errori di sintassi e di ortografia} & Errori di sintassi e di ortografia devono essere rimossi.\\
\hline
\textbf{Errori di battitura} & Possono essere presenti piccoli errori di battitura o di distrazione.\\
\hline
\textbf{Forma dei verbi} & I documenti devono essere scritti in forma verbale impersonale, e sarebbe inoltre preferibile l'utilizzo del presente indicativo.\\
\hline
\textbf{Backend e frontend} & Può essere utilizzata la forma con o senza ‘-’, pertanto si ritiene corretto sia l’uso di "back-end" e "front-end", sia di "backend" e "frontend".\\
\hline
\textbf{Formalità} & Qualsiasi espressione non formale deve essere sostituita da una formale.\\

\end{xltabular}
\caption{Punti di controllo per l'ispezione degli errori di forma e ortografici}
\end{table}



\subsubsection{Analisi dei Requisiti}

{\renewcommand{\arraystretch}{1.5}
\begin{table}[H]
\begin{xltabular}{\textwidth}{c|X}


\textbf{Titolo} & \textbf{Descrizione}   \\
\endfirsthead

\textbf{Titolo} & \textbf{Descrizione}   \\
\endhead

\multicolumn{2}{r}{{Continua nella pagina successiva}} \\
\endfoot

\endlastfoot

\hline
\textbf{Associazione UC - R}  &   Ogni caso d’uso dovrebbe avere un legame con uno o più requisiti.\\
\hline
\textbf{Ordinamento UC} & La numerazione dei casi d’uso dovrebbe rispecchiare la struttura di estensione.\\
\hline
\textbf{Diagrammi UML per UC} & (\textbf{da controllare}) Le estensioni di un caso d’uso dovrebbero essere rappresentate nello stesso diagramma UML del caso d’uso principale.\\
\hline
\textbf{Coerenza UC}  &   Verificare la coerenza tra la descrizione e il diagramma dei casi d'uso.\\
\hline
\textbf{Chiarezza Requisiti non funzionali} & I requisiti non funzionali dovrebbero essere definiti in modo chiaro.\\
\hline
\textbf{Completezza Tracciamento requisiti} & Il tracciamento dei requisiti dovrebbe essere completo e senza errori.\\
\hline
\textbf{Prioritizzazione Requisiti} & La priorità dei requisiti dovrebbe essere definita in modo coerente con le esigenze del progetto.\\

\end{xltabular}
\caption{Punti di controllo per l'ispezione del documento di Analisi dei Requisiti}
\end{table}

\subsubsection{Codifica}

{\renewcommand{\arraystretch}{1.5}
\begin{table}[H]
\begin{xltabular}{\textwidth}{c|X}


\textbf{Titolo} & \textbf{Descrizione}   \\
\endfirsthead

\textbf{Titolo} & \textbf{Descrizione}   \\
\endhead

\multicolumn{2}{r}{{Continua nella pagina successiva}} \\
\endfoot

\endlastfoot

\hline
\textbf{Nomi significativi} & Ogni componente del codice, come metodi e classi, dovrebbe avere un nome che ne descriva chiaramente la funzione.\\
\hline
\textbf{Case scorretto} & Le convenzioni descritte nelle Norme di progetto per i nomi delle classi o costanti devono essere rispettate.\\
\hline
\textbf{Commenti adeguati} & I metodi importanti dovrebbero essere preceduti da un commento che ne descriva lo scopo.\\
\hline
\textbf{Intestazioni mancanti} & Le intestazioni devono essere sempre presenti e conformi alle norme.\\
\hline
\textbf{Aggiornamento test} & Se un test produce un esito diverso da quanto riportato nel Piano di Qualifica, quest'ultimo deve essere aggiornato.\\

\end{xltabular}
\caption{Punti di controllo per l'ispezione del codice}
\end{table}





\section{Dashboard metriche impiegate}
\subsection{Misurazione e verifica} Sezione in cui vengono riportati i valori di misurazione delle metriche in uso. \color{red} Da implementare con strumenti e da misurare per ogni periodo a partire dal terzo o dal quarto periodo. \color{black}\\
\subsection{M1PMS}
\subsection{M2RNP}
\subsection{M3VP}
\subsection{M4VC}
\subsection{M5ISR}
\subsection{M6CCM}
\subsection{M7SC}
\subsection{M8BC}
\subsection{M9IG}
\subsubsection{Analisi dei requisiti}
\subsubsection{Norme di progetto}
\subsubsection{Piano di progetto}

\subsubsection{Piano di qualifica}
\subsection{M10PRI}
\subsection{M11LMC}
\subsection{M12DC}






\subsection{Processi primari}
    
   
\subsection{Processi di supporto}
\begin{comment}
    {\renewcommand{\arraystretch}{1.5}
    \begin{tabularx}{\textwidth}{p{0.05\textwidth}|p{0.35\textwidth}|X|X}
    \textbf{ID} & \textbf{Nome metrica} & \textbf{Valore accettabile} & \textbf{Valore ottimale}  \\
    \hline
    \multicolumn{4}{l}{\cellcolor{primarycolor}\textbf{\textit{Verifica}}} \\
    \hline
     &  &  &  \\
    \hline
     &  &  &  \\
    \hline
     &  &  &  \\
    \hline
     &  &  &  \\
    \hline
    \multicolumn{4}{l}{\cellcolor{primarycolor}\textbf{\textit{Gestione della qualità}}} \\
    \hline
     &  &  &  \\
    \end{tabularx}}  
  
\subsubsection{Documentazione}
    {\renewcommand{\arraystretch}{1.5}
    \begin{tabularx}{\textwidth}{p{0.15\textwidth}|p{0.35\textwidth}|X|X}
    \textbf{ID} & \textbf{Nome metrica} & \textbf{Valore accettabile} & \textbf{Valore ottimale}  \\
    \hline
    \multicolumn{4}{l}{\cellcolor{primarycolor}\textbf{\textit{Funzionalità}}} \\
    \hline
    MPD-CR & Copertura dei requisiti & 100 & 100 \\
    
    \hline
   
    \multicolumn{4}{l}{\cellcolor{primarycolor}\textbf{\textit{Usabilità}}}
    \\
    \hline
     MPD-DIU & Difficoltà incontrate dall'utente & 2  & 0 \\
      
    \hline
   
    \multicolumn{4}{l}{\cellcolor{primarycolor}\textbf{\textit{Efficienza}}}
    \\
    \hline
     MPD-TA & Tempo di apprendimento & 10 minuti & 5 minuti \\
      \hline
     \multicolumn{4}{l}{\cellcolor{primarycolor}\textbf{\textit{Affidabiltà}}} \\
     MPD-TM & Tempo di risposta medio & 3 secondi & 2 secondi \\
      \hline
     \multicolumn{4}{l}{\cellcolor{primarycolor}\textbf{\textit{Portabiltà}}} \\
     MPD-VS & La capacità del programma di eseguire su più piattaforme & 100  & 100 \\
    \hline
     \multicolumn{4}{l}{\cellcolor{primarycolor}\textbf{\textit{Manutenibilità}}} \\
     MPD-CC & la capacità del software di essere modificato, 
     includendo correzioni, miglioramenti o adattamenti. & 30  & 45 \\
    \end{tabularx}}   
\end{comment} 
\subsection{Processi organizzativi}

\section{Valutazioni per il miglioramento}  
\begin{comment}
\section{Resoconto delle attività di verifica}

\subsection{Verifica della documentazione}
\subsubsection{Errori ortografici}
\subsubsection{Indici di Gulpease}

\subsection{Verifica dei processi}

\subsubsection{Estimated at completion}
\subsubsection{Budget variance e schedule variance}
\subsubsection{Actual cost e estimate to complete}
\subsubsection{Earned value e planned value}
\subsubsection{Requirements stability index e satisfied obligatory requirements}
\subsubsection{Code coverage back-end}
\subsubsection{Code coverage front-end}
\subsubsection{Passed test cases percentage}
\subsubsection{Failed test cases percentage}
\subsubsection{Comprensibilità del codice}
\end{comment}
\end{document}
